\documentclass[11pt,a4paper]{article}
\usepackage[T1]{fontenc}
\usepackage{lmodern,amsmath,amssymb,bm,booktabs,array}
\usepackage[margin=25mm]{geometry}
\usepackage[unicode,colorlinks=true,linkcolor=blue,urlcolor=blue,citecolor=blue]{hyperref}
\usepackage{microtype}
\usepackage{fancyhdr}
\pagestyle{fancy}
\fancyhf{}
\fancyhead[L]{\small Censored gamma family: derivative derivation}
\fancyhead[R]{\small Draft for checking}
\fancyfoot[C]{\thepage}
\setlength{\headheight}{14pt}
\setlength{\parskip}{5pt}
\setlength{\parindent}{0pt}
\newcommand{\E}{\mathbb{E}}
\newcommand{\Var}{\operatorname{Var}}
\newcommand{\tr}{\operatorname{tr}}
\newcommand{\dd}{\mathrm{d}}
\newcommand{\code}[1]{\texttt{#1}}
\newcommand{\J}{\mathcal{J}}
\title{A censored gamma family for generalized additive models\\[5pt]
\large Likelihood, derivatives and prediction targets}
\author{Implementation note for \texttt{inlaws}\\
\small Prepared with Codex; draft for mathematical review}
\date{19 September 2026}
\begin{document}
\maketitle
\thispagestyle{plain}

This note accompanies \code{R/cgamma.R} on the \code{codex/cgamma} branch.
It follows the separation between fitting and prediction in Miller's censored
lognormal note \cite{miller}. The implemented family has one mean predictor,
a log link, and a single global dispersion parameter. It is not a gamma
location--scale model. All logarithms below are natural.

\section{Model and notation}
For independent responses, conditional on the model coefficients,
\begin{equation}
 Y_i\mid\bm\beta,\phi\sim\operatorname{Gamma}
 \left(k,\ \text{rate}=\frac{k}{\mu_i}\right),\qquad
 \eta_i=\bm x_i^\top\bm\beta+o_i,\quad \mu_i=e^{\eta_i},\quad k=\phi^{-1}.
 \label{eq:model}
\end{equation}
Thus $\E(Y_i\mid\bm\beta,\phi)=\mu_i$ and
$\Var(Y_i\mid\bm\beta,\phi)=\phi\mu_i^2$. Here $o_i$ is a known offset.
Internally the estimated family parameter is
\begin{equation}
 \theta=\log\phi,\qquad k=e^{-\theta},\qquad
 \partial_\theta k=-k,\quad \partial_\theta^2 k=k.
 \label{eq:theta}
\end{equation}
The public argument \code{theta} specifies dispersion on its original scale;
\code{getTheta()} returns $\theta$, whereas \code{getTheta(TRUE)} returns $\phi$.
This distinction is essential when checking derivatives.

Write $f(y;\mu,k)$ for the response density and $p_k(x)$ for the unit-rate gamma
density. Define the regularised tails
\[
 P(k,x)=\frac{1}{\Gamma(k)}\int_0^x t^{k-1}e^{-t}\dd t,
 \qquad Q(k,x)=1-P(k,x).
\]
For a fixed response bound $b$, use $x_b=kb/\mu$. In particular,
\[
 \partial_\mu x_b=-x_b/\mu,\qquad
 \partial_\theta x_b=-x_b,\qquad \partial_\theta^2x_b=x_b.
\]
Partial derivatives in $\mu$ hold $\theta$ fixed, and vice versa. The digamma
and trigamma functions are $\psi(k)$ and $\psi_1(k)$.

Prior weights $w_i\geq0$ multiply log likelihoods: $\ell=\sum_iw_i\ell_i$.
They are replication weights, not precision weights changing $k_i$.
The additional \code{mgcv} scale is fixed at one because $\phi$ is already
estimated as a family parameter.

\newpage
\section{Observed likelihood and saturated likelihood}
Drop the observation subscript until needed. Censoring is assumed
non-informative conditional on the covariates. With fixed censoring limits,
\begin{equation}
 \ell(\mu,\theta)=
 \begin{cases}
 \log f(y;\mu,k), & \text{exact }y>0,\\
 \log P(k,kc/\mu), & Y\leq c,\\
 \log Q(k,kc/\mu), & Y\geq c,\\
 \log\{P(k,ku/\mu)-P(k,kl/\mu)\}, & l\leq Y\leq u.
 \end{cases}
 \label{eq:lik}
\end{equation}
For a continuous response, endpoint inclusion does not affect probabilities.
One-sided censoring is an interval with $l=0$ or $u=\infty$.
Denote its probability by $A(\mu,\theta)$.

\subsection{A common deviance definition}
For each observation let
\[
 \ell_s(\theta)=\sup_{\mu>0}\ell(\mu,\theta),\qquad
 D(\mu,\theta)=2w\{\ell_s(\theta)-\ell(\mu,\theta)\}.
\]
For exact data $\mu_s=y$. For a proper one-sided censoring event the
supremum probability is one, so $\ell_s=0$. This also holds for
$[0,u]$, $[l,\infty)$ and $[0,\infty)$; the supremum need not be attained at
a finite positive mean.

For a finite interval $0<l<u<\infty$, define
\begin{equation}
 G(x,k)=xp_k(x)=\frac{x^ke^{-x}}{\Gamma(k)}.
 \label{eq:G}
\end{equation}
Differentiating the probability gives
\begin{equation}
 A_\mu=\frac{G(x_l,k)-G(x_u,k)}{\mu}.
 \label{eq:Amu}
\end{equation}
The stationary equation is therefore
\[
 0=\log G(x_l,k)-\log G(x_u,k)
   =k\log(l/u)+k(u-l)/\mu.
\]
It has the unique solution
\begin{equation}
 \boxed{\mu_s=\frac{u-l}{\log(u/l)}}.
 \label{eq:musat}
\end{equation}
The derivative is positive below $\mu_s$ and negative above it; interval
probability tends to zero at either end of the positive mean domain.
Thus this is the maximum. Crucially, $\mu_s$ is independent of $\theta$:
\[
 \ell_s^{(j)}(\theta)=\partial_\theta^j\log A(\mu_s,\theta),\qquad j=1,2.
\]
No derivative of an optimised mean is needed. Numerically evaluate
$\log(u/l)$ as $\log\{1+(u-l)/l\}$ using \code{log1p}.

The identity used to check likelihood normalisation is
\begin{equation}
 \frac12\sum_iD_i-\sum_iw_i\ell_{s,i}=-\sum_iw_i\ell_i.
 \label{eq:normalisation}
\end{equation}

\newpage
\section{Exact observations: all required derivatives}
Put $t=y/\mu$ and $x=kt$. The exact log density is
\begin{equation}
 \ell=k\log k-k\log\mu-\log\Gamma(k)+(k-1)\log y-ky/\mu.
 \label{eq:density}
\end{equation}
Since $\partial_\mu t=-t/\mu$, repeated differentiation gives
\begin{align}
 \ell_\mu&=k(t-1)/\mu,&
 \ell_{\mu\mu}&=k(1-2t)/\mu^2,\\
 \ell_{\mu\mu\mu}&=2k(3t-1)/\mu^3,&
 \ell_{\mu\mu\mu\mu}&=6k(1-4t)/\mu^4.
 \label{eq:exactmu}
\end{align}
Equivalently, for $r\geq1$,
\begin{equation}
 \ell_{\mu^r}=(-1)^r(r-1)!\,k(1-rt)\mu^{-r}.
 \label{eq:exactgeneral}
\end{equation}
Because $t$ is independent of $\theta$, each expression is proportional to
$k=e^{-\theta}$, and hence
\begin{equation}
 \boxed{\ell_{\mu^r\theta}=-\ell_{\mu^r},\qquad
        \ell_{\mu^r\theta\theta}=\ell_{\mu^r}.}
 \label{eq:exactmixed}
\end{equation}
Only $r\leq3$ is needed with one dispersion derivative, and $r\leq2$ with two.

\subsection{Pure dispersion derivatives}
First differentiate \eqref{eq:density} in $k$:
\[
 \ell_k=\log k+1-\psi(k)+\log t-t,
 \qquad \ell_{kk}=1/k-\psi_1(k).
\]
Applying $\partial_\theta=-k\partial_k$ yields
\begin{align}
 s:=\ell_\theta
   &=-k\{\log k+1-\psi(k)+\log t-t\},\label{eq:score}\\
 \ell_{\theta\theta}
   &=k\ell_k+k^2\ell_{kk}
    =-s+k-k^2\psi_1(k).\label{eq:score2}
\end{align}
Evaluate these at $t=1$ to obtain exact-observation saturated derivatives.
As a useful consistency check,
\[
 D=2wk\{t-1-\log t\}
\]
is the usual gamma deviance divided by dispersion. Consequently
$D_\theta=-D$ and $D_{\theta\theta}=D$ for exact observations.

\subsection{Independent integral interpretation}
For an interval with fixed bounds, differentiation under the integral gives
\begin{equation}
 A_\theta=\int_l^u f(y)s(y)\dd y,\qquad
 A_{\theta\theta}=\int_l^u f(y)
       \{s(y)^2+\ell_{\theta\theta}(y)\}\dd y.
 \label{eq:integrals}
\end{equation}
Thus $(\log A)_\theta=\E[s(Y)\mid l\leq Y\leq u]$ and
$(\log A)_{\theta\theta}=\E[\ell_{\theta\theta}(Y)\mid l\leq Y\leq u]
+\Var[s(Y)\mid l\leq Y\leq u]$. These identities provide an independent
quadrature check on the tail-derivative algorithm.

\newpage
\section{Censored observations: mean derivatives at the boundaries}
From \eqref{eq:G}, holding $k$ fixed,
\[
 \partial_\mu\log G(x_b,k)=(x_b-k)/\mu.
\]
For any polynomial $B(x,k)$ and integer $r\geq1$,
\begin{equation}
 \partial_\mu\{\mu^{-r}G(x,k)B(x,k)\}
 =\mu^{-r-1}G(x,k)\{(x-k-r)B-xB_x\}.
 \label{eq:recurrenceproof}
\end{equation}
Start from \eqref{eq:Amu} and set $B_0=1$. Induction gives
\begin{align}
 B_r(x,k)&=(x-k-r)B_{r-1}(x,k)-x\partial_xB_{r-1}(x,k),\label{eq:Brec}\\
 A_{\mu^r}&=\mu^{-r}\{G(x_l,k)B_{r-1}(x_l,k)
                     -G(x_u,k)B_{r-1}(x_u,k)\}.
 \label{eq:Ar}
\end{align}
The necessary polynomials are
\begin{align}
 B_0&=1,\\
 B_1&=x-k-1,\\
 B_2&=x^2-2(k+2)x+(k+1)(k+2),\\
 B_3&=x^3-3(k+3)x^2+3(k+2)(k+3)x
          -(k+1)(k+2)(k+3).
 \label{eq:Bexplicit}
\end{align}
For $k>0$, the boundary terms, and their required derivatives, vanish as
$x\downarrow0$ or $x\to\infty$. In computation they are set to zero at
these endpoints rather than evaluating $0\log0$ or $\infty-\infty$.

\subsection{Convert probability derivatives to log-probability derivatives}
Let $R_r=A_{\mu^r}/A$, and write $L_r=\partial_\mu^r\log A$. Repeated
application of the quotient rule gives
\begin{align}
 L_1&=R_1,\\
 L_2&=R_2-R_1^2,\\
 L_3&=R_3-3R_1R_2+2R_1^3,\\
 L_4&=R_4-4R_1R_3-3R_2^2+12R_1^2R_2-6R_1^4.
 \label{eq:logderivs}
\end{align}
For example, $\partial_\mu R_r=R_{r+1}-R_rR_1$, which verifies each line
from the previous one.

Compute the normalised boundary terms as
\[
 \frac{G(x_b,k)}{A}=\exp\{\log G(x_b,k)-\log A\}.
\]
This avoids forming a potentially underflowed probability $A$ explicitly.
The same expressions cover left, right and finite-interval censoring; only
the boundaries change. Narrow finite intervals require the alternative in
Section~\ref{sec:narrow} to avoid subtracting nearly identical boundary terms.

\newpage
\section{Mixed derivatives through second-order forward differentiation}
Represent a scalar function of $\theta$ by the triple
\[
 \J(a)=(a,\dot a,\ddot a),\qquad
 \dot a=\partial_\theta a,\quad\ddot a=\partial_\theta^2a.
\]
These are derivatives, not Taylor coefficients. Addition and subtraction are
componentwise; constants have zero derivative components. Product, reciprocal,
logarithm and exponential rules are
\begin{align}
 \J(ab)&=(ab,\dot a b+a\dot b,\ddot a b+2\dot a\dot b+a\ddot b),\\
 \J(a^{-1})&=\left(a^{-1},-\dot a/a^2,
                   2\dot a^2/a^3-\ddot a/a^2\right),\\
 \J(\log a)&=\left(\log a,\dot a/a,
                   \ddot a/a-(\dot a/a)^2\right),\\
 \J(e^a)&=e^a(1,\dot a,\ddot a+\dot a^2).
 \label{eq:jets}
\end{align}
All quotient operations below mean multiplication by the reciprocal triple.

\subsection{Boundary triples}
For $x=kb/\mu$, initialise $\J(k)=(k,-k,k)$ and $\J(x)=(x,-x,x)$.
Writing $q=\log G=k\log x-x-\log\Gamma(k)$ gives
\begin{align}
 \dot q&=x-k\{\log x+1-\psi(k)\},\label{eq:qdot}\\
 \ddot q&=-\dot q+k-k^2\psi_1(k).\label{eq:qddot}
\end{align}
Here both $x$ and $k$ vary with $\theta$. Treating the argument of the gamma
CDF as fixed when differentiating shape would omit part of the derivative.

Given $\J(\log A)$ from Sections~\ref{sec:series}--\ref{sec:cf}, form
$\J(G/A)=\J[\exp(q-\log A)]$. Evaluate the polynomials in
\eqref{eq:Bexplicit} using triples, construct $\J(R_r)$ using
\eqref{eq:Ar}, and apply \eqref{eq:logderivs} using the product rule above.
The output $\J(L_r)$ contains $\ell_{\mu^r}$,
$\ell_{\mu^r\theta}$ and $\ell_{\mu^r\theta\theta}$.

\subsection{Explicit mixed expressions for checking}
The first two orders, written without triple notation, are
\begin{align}
 \dot L_1&=\dot R_1,& \ddot L_1&=\ddot R_1,\\
 \dot L_2&=\dot R_2-2R_1\dot R_1,&
 \ddot L_2&=\ddot R_2-2\{\dot R_1^2+R_1\ddot R_1\},\\
 \dot L_3&=\dot R_3-3\{\dot R_1R_2+R_1\dot R_2\}
                +6R_1^2\dot R_1.\nonumber
\end{align}
Only these mixed expressions are required by \code{Dd}; $L_4$ itself needs no
dispersion derivative. The implementation carries a complete triple through
the polynomial arithmetic for simplicity. No finite difference is used here.

\newpage
\section{Dispersion derivatives of gamma tails: lower-tail series}
\label{sec:series}
The incomplete-gamma expansion in \cite{dlmfseries} can be written
\begin{equation}
 P(k,x)=\frac{e^{q(x,k)}}{k}\sum_{j=0}^{\infty}T_j,
 \quad T_0=1,\qquad T_j=T_{j-1}\frac{x}{k+j}.
 \label{eq:series}
\end{equation}
The implementation uses this expansion for $x<k+1$. All terms are positive.
At fixed observation bound and mean, $\dot x=-x$ and $\dot k=-k$.
For the ratio $r_j=x/(k+j)$,
\[
 \partial_\theta\log r_j=-1+\frac{k}{k+j}=-\frac{j}{k+j},
 \qquad
 \partial_\theta^2\log r_j=-\frac{jk}{(k+j)^2}.
\]
Consequently an explicit derivative recurrence is
\begin{align}
 \dot r_j&=-r_j\frac{j}{k+j},\\
 \ddot r_j&=r_j\frac{j(j-k)}{(k+j)^2},\\
 \dot T_j&=\dot T_{j-1}r_j+T_{j-1}\dot r_j,\\
 \ddot T_j&=\ddot T_{j-1}r_j+2\dot T_{j-1}\dot r_j+T_{j-1}\ddot r_j.
\end{align}
Initialise $\dot T_0=\ddot T_0=0$. Sum these sequences to obtain
$S=\sum T_j$, $\dot S=\sum\dot T_j$, $\ddot S=\sum\ddot T_j$.
Taking the logarithm in \eqref{eq:series} gives
\begin{align}
 \log P&=q-\log k+\log S,\\
 \partial_\theta\log P&=\dot q+1+\dot S/S,\\
 \partial_\theta^2\log P&=\ddot q+\ddot S/S-(\dot S/S)^2.
 \label{eq:seriesderivs}
\end{align}
The $+1$ in the first derivative is from $-\log k=\theta$; there is no
corresponding second-derivative term.

\subsection{Implementation and convergence}
The code computes the recurrence with the rules in \eqref{eq:jets} rather
than maintaining separate hand-expanded recurrences. It stops when all three
components of the new term satisfy
\[
 |\J(T_j)_m|<2\times10^{-14}\{1+|\J(S)_m|\},\quad m=0,1,2.
\]
There is a hard limit of 10,000 iterations with an error on nonconvergence.
Convergence of the probability value alone is not sufficient for convergence
of its derivatives.

R's \code{pgamma(log.p = TRUE)} supplies the final log-probability value.
The differentiated series supplies its first two dispersion derivatives.
This preserves R's specialised tail evaluation without numerically
perturbing the dispersion. Very large shapes near the transition region
remain a regime to test beyond the present finite validation grid.

\newpage
\section{Dispersion derivatives of gamma tails: upper-tail fraction}
\label{sec:cf}
An equivalent contracted form of the upper incomplete-gamma continued
fraction \cite{dlmfcf} is
\begin{equation}
 Q(k,x)=e^{q(x,k)}H(k,x),\quad
 H=\cfrac{1}{b_0+\cfrac{a_1}{b_1+\cfrac{a_2}{b_2+\cdots}}},
 \quad b_j=x+1-k+2j,\quad a_j=j(k-j).
 \label{eq:cf}
\end{equation}
Use this representation for $x\geq k+1$. The triple inputs for its
coefficients are particularly simple:
\begin{equation}
 \J(b_j)=(b_j,-x+k,x-k),\qquad
 \J(a_j)=(a_j,-jk,jk).
 \label{eq:cfinputs}
\end{equation}

\subsection{Differentiating the fraction evaluation}
A multiplicative evaluation starts with
\[
 C_0=M,\quad D_0=b_0^{-1},\quad H_0=D_0,\qquad M=10^{300},
\]
where $M$ approximates an infinite initial $C_0$ and its derivative components
are zero. For $j=1,2,\ldots$ compute
\begin{align}
 D_j&=(b_j+a_jD_{j-1})^{-1},\\
 C_j&=b_j+a_j/C_{j-1},\\
 \Delta_j&=C_jD_j,\\
 H_j&=H_{j-1}\Delta_j.
 \label{eq:cfalgorithm}
\end{align}
Replace each operation by its triple rule in \eqref{eq:jets}. This yields
$H$, $\dot H$ and $\ddot H$, and hence
\begin{align}
 \log Q&=q+\log H,\\
 \partial_\theta\log Q&=\dot q+\dot H/H,\\
 \partial_\theta^2\log Q&=\ddot q+\ddot H/H-(\dot H/H)^2.
 \label{eq:cfderivs}
\end{align}
The stopping rule checks $|\Delta_j-1|$, $|\dot\Delta_j|$ and
$|\ddot\Delta_j|$ against $2\times10^{-14}$, with the same iteration limit
as the series. Derivatives of a fraction that terminates in value at integer
$k$ can still require further iterations: the derivative tests must remain.

The upper-tail value is again supplied by \code{pgamma}, with
\code{lower.tail = FALSE}. Derivatives are obtained from the differentiated
fraction. A requested tail opposite to the directly evaluated tail is
obtained by the log-difference rule on the next page.

These expressions specify a numerical evaluation of analytic derivatives.
They are not finite differences, but finite iteration and floating-point
arithmetic still require numerical validation. No assertion of uniform
accuracy over all positive shapes is made by the present implementation.

\newpage
\section{Stable interval probabilities and narrow intervals}
\label{sec:narrow}
Let $a=\log A_1$, $b=\log A_0$ with $A_1>A_0\geq0$ and define
$z=\log(A_1-A_0)$. Evaluate
\begin{equation}
 z=a+\log\{-\operatorname{expm1}(b-a)\},\quad
 r=e^{b-z},\quad d=\dot a-\dot b.
\end{equation}
Direct differentiation gives
\begin{align}
 \dot z&=\dot a+rd,\\
 \ddot z&=\ddot a+r(\ddot a-\ddot b)-r(1+r)d^2.
 \label{eq:logdiff}
\end{align}
This also computes a complement by setting $a=0$ with zero derivatives.
At a tail endpoint having probability zero, use log value $-\infty$ and
zero derivative components, so $r=0$.

For a finite interval either subtract lower tails,
$P(k,x_u)-P(k,x_l)$, or upper tails, $Q(k,x_l)-Q(k,x_u)$.
The current implementation uses lower tails if $P(k,x_l)<1/2$, otherwise
upper tails. This avoids subtracting two near-one probabilities in the far
upper tail, although it cannot remove all cancellation for narrow intervals.

\subsection{Integrate the density and its derivatives together}
For $\delta=u-l$ small, substitute $y=l+\delta(v+1)/2$:
\begin{equation}
 A=\frac{\delta}{2}\int_{-1}^{1}f\{l+\delta(v+1)/2\}\dd v
   \simeq\frac{\delta}{2}\sum_{j=1}^{16}\omega_j f(y_j).
 \label{eq:quad}
\end{equation}
The code uses 16-point Gauss--Legendre quadrature when
$0<l<u<\infty$, $\delta/l<10^{-4}$ and $k\delta/\mu<1$.
This is a local numerical approximation, not a replacement likelihood model.

For density derivatives define $C_0=1$ and
\[
 \partial_\mu^rf=f\mu^{-r}C_r(x,k),\qquad
 C_{r+1}=(x-k-r)C_r-x\partial_x C_r.
\]
Thus $C_1=x-k$, $C_2=x^2-2(k+1)x+k(k+1)$, and
\begin{align*}
 C_3&=x^3-3(k+2)x^2+3(k+1)(k+2)x-k(k+1)(k+2),\\
 C_4&=x^4-4(k+3)x^3+6(k+2)(k+3)x^2\\
    &\hspace{9mm}-4(k+1)(k+2)(k+3)x+k(k+1)(k+2)(k+3).
\end{align*}
Integrate $f C_r/\mu^r$ with the same nodes and weights, divide by the
integral of $f$, and use \eqref{eq:logderivs}. Carry all operations in
$\theta$ triples. A common subtraction of the largest node log density
prevents underflow; treat this numerical scaling constant as fixed within
the derivative calculation. As $\delta\downarrow0$, log probability approaches
$\log f\{(l+u)/2\}+\log\delta$ and its parameter derivatives approach those
of an exact observation. This is a separate regression test.

\newpage
\section{Avoiding cancellation in far upper-tail derivatives}
The boundary formula is algebraically correct but can be numerically poor
far into the upper tail. For example, $R_4$ and $R_1^4$ can both have order
$x^4$ while their combination in $L_4$ has order $x$. Subtracting these
large quantities can destroy the fourth derivative. Working on the log
probability scale alone does not resolve this cancellation.

Repeated integration by parts gives the upper-tail expansion
\begin{equation}
 Q(k,x)\sim \frac{G(x,k)}{x}S(x,k),\qquad
 S=\sum_{j\geq0}U_j,\quad U_0=1,\quad
 U_j=U_{j-1}\frac{k-j}{x}.
 \label{eq:upperexpansion}
\end{equation}
This is an asymptotic expansion; it must be truncated while the terms are
small, not summed indefinitely. The code uses it only for
$x>\max(100,5k)$ and requires convergence of all retained derivatives.
Since $x=kb/\mu$, $U_j$ is proportional to $\mu^j$ when $k$ is fixed:
\[
 S_{\mu^r}=\mu^{-r}\sum_{j\geq r}
 j(j-1)\cdots(j-r+1)U_j,\qquad r=1,\ldots,4.
\]
Evaluate these sums as $\theta$ triples, normalise by $S$, and use
\eqref{eq:logderivs} to get derivatives of $\log S$. Then
\begin{equation}
 (\log Q)_{\mu^r}
 =(-1)^r(r-1)!\,\mu^{-r}\{k-1-rx\}+(\log S)_{\mu^r}.
 \label{eq:farmean}
\end{equation}
The dominant order-$x$ term is now explicit, and the remaining small
series does not require cancellation of order-$x^4$ terms. Dispersion and
mixed derivatives follow by carrying the same triples throughout.

\subsection{Finite intervals wholly in the far upper tail}
Write $L=\log Q(k,x_l)$, $U=\log Q(k,x_u)$, $\Delta=U-L$,
$e=\exp\Delta$ and $F=1-e$. Then $\log A=L+\log F$.
For $\Delta_r=\partial_\mu^r\Delta$,
\begin{align*}
 e_\mu&=e\Delta_1,\\
 e_{\mu^2}&=e(\Delta_2+\Delta_1^2),\\
 e_{\mu^3}&=e(\Delta_3+3\Delta_1\Delta_2+\Delta_1^3),\\
 e_{\mu^4}&=e(\Delta_4+4\Delta_1\Delta_3+3\Delta_2^2
                      +6\Delta_1^2\Delta_2+\Delta_1^4).
\end{align*}
Because $F_{\mu^r}=-e_{\mu^r}$, applying \eqref{eq:logderivs} to the ratios
$-e_{\mu^r}/F$ supplies $(\log F)_{\mu^r}$. Add these to the derivatives of
$L$ from \eqref{eq:farmean}. This calculation also uses triples for mixed
derivatives. The narrow-interval quadrature takes precedence when its
conditions hold. For $u=\infty$, $\log A=L$ directly.

\newpage
\section{Mapping the derivatives to the \texttt{mgcv} interface}
The extended-family interface requires derivatives of deviance in $\mu$,
not derivatives in $\eta$. For $r\geq1$ the saturated likelihood does not
depend on the fitted mean, giving
\[
 D_{\mu^r\theta^s}=-2w\ell_{\mu^r\theta^s},\qquad
 D_\theta=2w(\ell_{s,\theta}-\ell_\theta),\quad
 D_{\theta\theta}=2w(\ell_{s,\theta\theta}-\ell_{\theta\theta}).
\]
The return fields of \code{Dd(y, mu, theta, wt, level)} are:
\begin{center}
\begin{tabular}{lll}
\toprule
Level & Field & Derivative\\
\midrule
0 & \code{Dmu}, \code{Dmu2} & $D_\mu$, $D_{\mu\mu}$\\
1 & \code{Dmu3}, \code{Dth} & $D_{\mu\mu\mu}$, $D_\theta$\\
1 & \code{Dmuth}, \code{Dmu2th} & $D_{\mu\theta}$, $D_{\mu\mu\theta}$\\
2 & \code{Dmu4}, \code{Dth2} & $D_{\mu\mu\mu\mu}$, $D_{\theta\theta}$\\
2 & \code{Dmuth2}, \code{Dmu2th2} & $D_{\mu\theta\theta}$, $D_{\mu\mu\theta\theta}$\\
2 & \code{Dmu3th} & $D_{\mu\mu\mu\theta}$\\
\bottomrule
\end{tabular}
\end{center}
Levels are cumulative. For exact data use the Gamma expected information:
$\code{EDmu2}=2wk/\mu^2$ and $\code{EDmu2th}=-2wk/\mu^2$.
For censored data use the working approximation
$\code{EDmu2}=\max(0,D_{\mu\mu}+D_\mu/\mu)$ and, where positive,
$\code{EDmu2th}=D_{\mu\mu\theta}+D_{\mu\theta}/\mu$ (zero otherwise).
The latter uses non-negative observed curvature in $\eta$, not exact
expected information. Using $D_{\mu\mu}$ alone can give negative working
weights in \code{bam}, even when $D_{\eta\eta}>0$. Only the exact-data
choice is required to recover Gamma's EDF and covariance conventions.

The \code{ls} function returns $\sum_iw_i\ell_{s,i}$, its first two
$\theta$ derivatives, and the per-observation first derivatives as
\code{LSTH1}. The \code{aic} function returns $-2\sum_iw_i\ell_i$;
\code{mgcv} adds the model complexity contribution. Zero weights contribute
zero to all of these quantities.

\subsection{Log-link conversion and checks}
Since $\partial_\eta=\mu\partial_\mu$,
\begin{align}
 D_\eta&=\mu D_\mu,\\
 D_{\eta\eta}&=\mu^2D_{\mu\mu}+\mu D_\mu,\\
 D_{\eta^3}&=\mu^3D_{\mu^3}+3\mu^2D_{\mu^2}+\mu D_\mu,\\
 D_{\eta^4}&=\mu^4D_{\mu^4}+6\mu^3D_{\mu^3}+7\mu^2D_{\mu^2}+\mu D_\mu.
 \label{eq:link}
\end{align}
The same linear combinations hold after a $\theta$ derivative because
$\mu$ is fixed in that derivative. The family supplies the $\mu$ derivatives;
\code{mgcv} performs this conversion.

For an exact observation, direct differentiation of
$\ell=\text{constant}-k\eta-ky e^{-\eta}$ gives
$\ell_\eta=x-k$, $\ell_{\eta^2}=-x$, $\ell_{\eta^3}=x$ and
$\ell_{\eta^4}=-x$. These independently check \eqref{eq:link}.
For uncensored data the fitted model matches an ordinary Gamma fit under
ML or REML, including likelihood dispersion, coefficients and EDF. The
extended-family smoothing parameter equals the Gamma smoothing parameter
divided by $\phi$. Gamma stores its likelihood dispersion in
\code{reml.scale}; \code{sig2} and \code{scale} use a separate Fletcher
estimate by default. To compare covariance matrices and standard errors,
rescale the ordinary Gamma covariance by \code{reml.scale / sig2}.
The censored-family external scale stays at one, with dispersion in
\code{family\$getTheta(TRUE)}.

\newpage
\section{Prediction: conditional mean and epsilon correction}
\label{sec:prediction}
The implemented response prediction is
\[
 \widehat\mu_0=g(\widehat{\bm\beta})
  =\exp(\bm x_0^\top\widehat{\bm\beta}+o_0).
\]
This is already a plug-in estimate of the conditional arithmetic mean.
There is no lognormal residual-variance adjustment. A different target is
$\E\{g(\bm\beta)\mid\bm y,\widehat\phi,\widehat{\bm\lambda}\}$,
which averages a nonlinear function over coefficient uncertainty.
The epsilon approach of Thorson and Kristensen \cite{tk}, based on Tierney,
Kass and Kadane \cite{tierney}, is relevant to this target.

\subsection{Expansion at the penalised mode}
Fix dispersion and smoothing parameters for this derivation. Define the
negative log posterior, up to a constant, by
\[
 h(\bm\beta)=-\sum_iw_i\ell_i(\eta_i,\theta)
            +\tfrac12\bm\beta^\top S\bm\beta.
\]
At its mode let $V=\{h''(\widehat{\bm\beta})\}^{-1}$,
$T=h'''(\widehat{\bm\beta})$ and $Q=h''''(\widehat{\bm\beta})$.
Unlike Miller's per-observation $h$ convention, $h$ here is the full
negative log posterior, so there are no explicit factors of sample size.
The fully exponential Laplace expansion gives
\begin{equation}
 \E(g\mid\bm y)\simeq
 g+\tfrac12\tr(Vg'')-\tfrac12 T[V,Vg'].
 \label{eq:epsilon}
\end{equation}
In components $T[V,Vg']=\sum_{abc}T_{abc}V_{ab}(Vg')_c$.
The last term corrects for departure from a Gaussian posterior and must not
be dropped merely because a Gaussian covariance estimate is available.

For $g=e^{\bm x_0^\top\bm\beta+o_0}$,
$g'=g\bm x_0$ and $g''=g\bm x_0\bm x_0^\top$. Also
\[
 T=\sum_it_i\bm x_i^{\otimes3},\quad
 Q=\sum_iq_i\bm x_i^{\otimes4},\quad
 t_i=\tfrac12D_{i,\eta^3},\quad q_i=\tfrac12D_{i,\eta^4}.
\]
The quadratic penalty has no third or fourth derivative. Set
$v_0=\bm x_0^\top V\bm x_0$,
$d_i=\bm x_i^\top V\bm x_i$ and $e_i=\bm x_i^\top V\bm x_0$.
Then \eqref{eq:epsilon} reduces to
\begin{equation}
 \boxed{\E(g\mid\bm y)\simeq
 g\left\{1+\tfrac12v_0-\tfrac12\sum_it_id_ie_i\right\}.}
 \label{eq:epsiloncontract}
\end{equation}
For exact gamma observations $t_i=-w_i k y_i/\mu_i$ and
$q_i=w_i k y_i/\mu_i$. Therefore \emph{all observations}, not just censored
ones, contribute to the correction. This prevents direct reuse of the
censored-row-only calculation for the lognormal family.

\newpage
\section{Checking the prediction expansion and its variance}
\subsection{An exact intercept-only check}
Consider $n$ uncensored observations, unit weights, fixed $k$, and a flat
prior on $\eta=\log\mu$. Write $a=nk$ and $b=k\sum_i y_i$.
The likelihood in $\eta$ is proportional to $\exp(-a\eta-be^{-\eta})$.
The change of variables $\mu=e^\eta$ shows
\[
 \mu\mid\bm y\sim\operatorname{InverseGamma}(a,b),\qquad
 \E(\mu\mid\bm y)=\frac{b}{a-1}\quad(a>1).
\]
At the mode in $\eta$, $g=b/a=\bar y$, $V=1/a$, $T=-a$ and $Q=a$.
Equation~\eqref{eq:epsilon} therefore yields
\[
 g+\frac{g}{2a}+\frac{g}{2a}=g(1+1/a),
\]
which agrees with the expansion of $b/(a-1)$ to this order. Retaining only
the curvature term gives $g(1+1/(2a))$ and misses half of the leading
correction. The exact posterior exists for $a>0$, but its mean requires
$a>1$; a local expansion cannot guarantee existence of posterior moments.

\subsection{Variance expression to be checked before implementation}
For a scalar prediction define
\[
 \bm a=Vg',\quad B=g''-T[\bm a],\quad
 \bm b=V\{2g''\bm a-T[\bm a,\bm a]\}.
\]
Here $T[\bm a]$ is a matrix, $T[\bm a,\bm a]$ a vector, and
$Q[\bm a,\bm a]$ a matrix. Differentiating the log determinant of the
perturbed Hessian in the fully exponential approximation gives
\begin{equation}
 \Var(g\mid\bm y)\simeq
 g'^\top Vg'
 +\tfrac12\tr\{(VB)^2\}
 -\tfrac12\tr\left[V\{Q[\bm a,\bm a]+T[\bm b]-2g'''[\bm a]\}\right].
 \label{eq:var}
\end{equation}
For the intercept-only check above, this equals
$g^2\{a^{-1}+4a^{-2}\}$, matching the expansion of the exact variance
$b^2/\{(a-1)^2(a-2)\}$ for $a>2$.

\subsection{Scope of the prediction result}
Equations~\eqref{eq:epsiloncontract} and \eqref{eq:var} are derivations for a
future prediction option, not a currently implemented prediction method.
They condition on estimated dispersion and smoothing parameters. Joint
integration over these quantities requires corresponding joint derivatives
and covariance information. Prediction for a \emph{new} random-effect level
also introduces a different integral and is not supplied by these formulae
alone.

The code presently follows ordinary \code{Gamma(link = "log")} prediction
semantics. Posterior-mean prediction, frequentist bias correction of fixed
parameters, and marginal prediction over new random effects should not be
presented as interchangeable targets. Numerical integration of small models
and simulation should precede exposing an epsilon-corrected option.

\newpage
\section{Verification and reviewer guide}
The companion regression checks are in \code{tests/cgamma.R}. They cover:
\begin{enumerate}
 \item First through fourth mean derivatives and every required mixed
 derivative, compared with Richardson differences of lower-order quantities.
 Exact, left-, right-, ordinary interval and narrow interval responses are
 included, together with support endpoints and zero weights.
 \item Pure dispersion derivatives of log tails, compared with independently
 perturbed calls to base R's \code{pgamma}. The shape grid is
 $0.03,0.2,1,2,10,100,1000$ and the bound-to-mean ratio grid is
 $0.001,0.1,0.9,1,1.1,2,10,100$, for both tails. Independent numerical integrals of the density score and its
 derivative also check dispersion derivatives for each censoring direction.
 Far upper-tail mean and mixed derivatives are checked separately.
 \item Saturated-likelihood derivatives and the exact likelihood identity
 \eqref{eq:normalisation}; gamma density and deviance formulae in the
 uncensored case.
 \item The narrow-interval limit, including derivatives; an ordinary gamma
 smooth fit at fixed and estimated dispersion (ML and REML, vector and
 two-column exact responses); and a censored regression fit against
 independent direct likelihood optimisation.
 \item Estimated-dispersion smooth fits, offsets, response predictions,
 preservation of censoring when subsetting, optimised null deviance, both
 \code{bam} fitting paths, and input validation.
\end{enumerate}
These checks detect algebraic and implementation errors; they are not a
proof of uniform accuracy or a large simulation study of estimator behaviour.
Tail calculations at more extreme shapes, censoring-driven weak
identifiability, and epsilon-corrected predictions deserve further study.

For review, the most consequential checks are the sign change
$\theta=\log\phi=-\log k$; the simultaneous change of $x$ and $k$ in
\eqref{eq:qdot}; the saturated interval mean in \eqref{eq:musat}; the
boundary recurrence \eqref{eq:Brec}; and the distinction between posterior
mean and plug-in prediction in Section~\ref{sec:prediction}.

\begin{thebibliography}{9}\small
\bibitem{miller} Miller, D. L. (2026).
\emph{Modelling censored log-normal responses in generalized additive models}.
28 August 2026. \code{inlaws}, \code{inst/maths/clogn/clogn.pdf}.
Especially Section 3 and Appendix equation (9).
\bibitem{dlmfseries} NIST Digital Library of Mathematical Functions.
Section 8.7, especially equation 8.7.1.
\url{https://dlmf.nist.gov/8.7}.
\bibitem{dlmfcf} NIST Digital Library of Mathematical Functions.
Section 8.9, continued fractions for incomplete gamma functions.
\url{https://dlmf.nist.gov/8.9}.
\bibitem{tk} Thorson, J. T. and Kristensen, K. (2016).
Implementing a generic method for bias correction in statistical models using
random effects, with spatial and population dynamics examples.
\emph{Fisheries Research}, 175, 66--74.
\href{https://doi.org/10.1016/j.fishres.2015.11.016}{doi:10.1016/j.fishres.2015.11.016}.
\bibitem{tierney} Tierney, L., Kass, R. E. and Kadane, J. B. (1989).
Fully exponential Laplace approximations to expectations and variances of
nonpositive functions. \emph{JASA}, 84(407), 710--716.
\bibitem{wood} Wood, S. N., Pya, N. and S\"afken, B. (2016).
Smoothing parameter and model selection for general smooth models.
\emph{JASA}, 111(516), 1548--1563.
\href{https://doi.org/10.1080/01621459.2016.1180986}{doi:10.1080/01621459.2016.1180986}.
\end{thebibliography}
\end{document}
